% BCT Chemistry Series — Volume K
% Master LaTeX document
% Letters 78–85 and Appendices KA1–KC2
% Michel Robert Cabrié — 20 March 2026
% ORCID: 0009-0007-9561-9859
%
% This master file collects the BCT Chemistry Series for submission
% to Zenodo as a single citable volume. Each letter and appendix
% is provided as a standalone section, formatted for the revtex4-2
% document class used throughout the BCT programme.
%
% To compile: pdflatex BCT_ChemistrySeries_VolK.tex (twice for TOC)
% Requires: revtex4-2, amsmath, booktabs, xcolor, hyperref

\documentclass[aps,preprint,superscriptaddress,floatfix]{revtex4-2}
\usepackage{amsmath,amssymb,amsfonts}
\usepackage{booktabs,array,xcolor,hyperref}
\usepackage{graphicx}

\definecolor{navy}{HTML}{0B2545}
\definecolor{teal}{HTML}{006064}
\definecolor{dgrn}{HTML}{1B5E20}

% ── VOLUME METADATA ──────────────────────────────────────────
\begin{document}

\title{BCT Chemistry Series --- Volume K:\\
Letters 78--85 and Appendices KA1--KC2}

\author{Michel Robert Cabri\'{e}}
\affiliation{Independent Researcher, Barrys Reef, Victoria, Australia}
\affiliation{ORCID: 0009-0007-9561-9859}

\date{20 March 2026}

\begin{abstract}
This volume collects the BCT Chemistry Series: eight Letters
(78--85) and nine companion Appendices (KA1--KC2), documenting
the extension of the BCT Superfluid Lattice Model from particle
physics into atomic structure, molecular geometry, and chemical
periodicity.

Key results: shell capacity $2n^2$ as an exact derived theorem
(not postulated); Pauli exclusion and noble gas inertness as
topological theorems; the algebraic identity
$r_\mathrm{tet}(1+r_\mathrm{tet}) = 1/8$ (exact);
H$_2$O bond angle to $-0.0001\%$; NH$_3$ bond angle to
$-0.027\%$ (after electronegativity correction, Letter 85);
carbon uniqueness ($S_H = 1$) as the topological basis of life;
transition metal anomalies as Hopf symmetry attractors;
a complete topological reclassification of all 118 elements;
and electronegativity from void geometry with P at $+0.1\%$
and O at $+0.2\%$.

All results from zero free parameters. BCT programme total:
122+ predictions, 41 sub-0.1\%, 0 free parameters,
3 geometric inputs.
\end{abstract}

\maketitle
\tableofcontents

% ════════════════════════════════════════════════════════════
% PART I: LETTERS
% ════════════════════════════════════════════════════════════

\newpage
\part*{Part I: Letters 78--85}

% Each letter is a standalone revtex document.
% For the combined volume, we include them as sections.
% The individual .tex files are provided separately for
% standalone compilation.

% ── LETTER 78 ────────────────────────────────────────────────
\newpage
\section*{Letter 78: Atomic Structure from OHC Bessel Resonance Modes}
\addcontentsline{toc}{section}{Letter 78: Atomic Structure}
% \documentclass[aps,prl,twocolumn,showpacs,superscriptaddress]{revtex4-2}
\usepackage{amsmath,amssymb,booktabs,xcolor,hyperref}

\begin{document}

\begin{abstract}
We derive atomic structure, the Pauli exclusion principle, the octet 
rule, and the chemical inertness of noble gases from BCT vacuum topology, 
without additional postulates. Electrons are H=1 Hopf defects of the 
Octet-Hopfion Condensate (OHC); protons are D4 instanton nodes; electron 
shells are quantised OHC Bessel resonance modes with capacity $2n^2$. 
The Pauli exclusion principle is the topological complementarity of 
$H=+1$ and $H=-1$ Hopfions, which combine to net $H=0$: a theorem, 
not a postulate. Noble gases are \textbf{closed Hopf multiplets} --- 
topologically inert, not merely energetically stable. The octet rule 
derives from the two BCT void types: $r_{\rm oct}$ generates 1 spherical 
(s) mode; $r_{\rm tet}$ generates 3 directional (p) modes; giving 
$4 \times 2 = 8$ electrons at $n=2$. Gilbert Lewis drew $r_{\rm oct}$ 
and $r_{\rm tet}$ in 1916. He called them dots. Chemistry is topology. 
The periodic table is a Hopf fibration atlas.
\end{abstract}

\maketitle

\section{Introduction}

Standard quantum mechanics derives atomic structure correctly but 
incompletely. It provides rules without reasons: the Pauli exclusion 
principle is postulated, not derived; the octet rule is a pattern, 
not a theorem; noble gas inertness is energetic, not topological. 
The BCT Superfluid Lattice Model provides the missing depth.

The recent result of White et al.~\cite{White2026} --- deriving the 
hydrogen spectrum acoustically from a dynamic vacuum via the 
Madelung/Bogoliubov framework --- is the limiting case of BCT. 
BCT specifies the exact vacuum geometry ($c/a = \sqrt{2}$, the unique 
Lorentz-invariant BCT lattice) and derives all coupling constants.

\section{BCT Atomic Constituents}

\begin{center}
\begin{tabular}{lll}
\toprule
Entity & BCT identity & Origin \\
\midrule
Proton/neutron & D4 instanton node & Letter~71 \\
Electron & $H=1$ Hopfion & $\pi_3(S^2)=\mathbb{Z}$ \\
Electron spin & OHC phase $\pm$ & Opposite orientation \\
Shell $n$ & OHC Bessel mode $n$ & Josephson at $\alpha_0$ \\
\bottomrule
\end{tabular}
\end{center}

\section{Hydrogen}

The hydrogen atom: single D4 node (proton) with single $H=1$ Hopfion 
(electron) in the first OHC Bessel resonance. The Bohr radius:
\begin{equation}
    a_0 = \frac{\hbar}{m_e c\,\alpha} = 0.5293~\text{\AA}
    \quad \text{(observed: } 0.5292~\text{\AA}\ \checkmark\text{)}
\end{equation}
where $\alpha = \alpha_0(1-2\alpha_0) = 1/137.018$ (Letter~1, error 
$-0.013\%$). Shell energies $E_n = -13.606/n^2~\text{eV}$ follow from 
the OHC Bessel mode spectrum; the $1/n^2$ ladder is exact.

\section{Pauli Exclusion as Topological Complementarity}

\begin{theorem}[BCT Pauli Theorem]
Two $H=+1$ Hopfions in the same spatial OHC resonance mode must have 
opposite OHC interior phase orientations (opposite spin). Two maps 
$S^3\to S^2$ with the \emph{same} orientation combine to $H=+2$, 
exceeding the single-mode capacity. Two maps with \emph{opposite} 
orientation ($H=+1$ and $H=-1$) combine to net $H=0$ --- topologically 
complementary --- and share the same mode.
\end{theorem}

The Pauli principle is not postulated. It is topological necessity: 
$H=+2$ exceeds the $n=1$ mode capacity. $H\in\mathbb{Z}$ follows from 
$\pi_3(S^2) = \mathbb{Z}$ (Appendix~JH).

\section{The Octet Rule}

The BCT vacuum has exactly two void types. Their geometry directly 
determines the orbital symmetries at $n=2$:
\begin{itemize}
\item $r_{\rm oct}$: 1 isotropic centre $\to$ 1 s-orbital $\to$ 2 electrons
\item $r_{\rm tet}$: 3 axial orientations $\to$ 3 p-orbitals $\to$ 6 electrons
\end{itemize}
Total: $4 \times 2 = \mathbf{8}$ electrons. The octet is the BCT void 
geometry at $n=2$. Lewis drew $r_{\rm oct} + 3r_{\rm tet}$ in 1916. 
He called them dots.

\section{Noble Gases as Closed Hopf Multiplets}

\begin{theorem}[BCT Noble Gas Theorem]
A noble gas is an atom whose outermost electron shell is a \textbf{closed 
Hopf multiplet}: all $2n^2$ Hopf modes at level $n$ are occupied, 
with net shell Hopf charge $= 0$.

A closed Hopf multiplet is \textbf{topologically inert}: it cannot 
accept an additional electron (which would require opening a new Hopf 
mode at the next resonance level) and cannot donate an electron 
(which would break the topological closure) without supplying 
energy $\propto \Delta E_{\rm top}$.
\end{theorem}

Noble gas inertness is not merely energetic. It is topological. 
The Hopf multiplet is complete. There is no available topological 
channel for bonding.

\begin{center}
\begin{tabular}{lllr}
\toprule
Gas & $Z$ & Outer config & BCT status \\
\midrule
He & 2 & $1s^2$ & $H{=}+1,H{=}-1$ pair, net $H{=}0$ \\
Ne & 10 & $2s^22p^6$ & oct+3 tet all filled \\
Ar & 18 & $3s^23p^6$ & $n=3$ valence closed \\
Kr & 36 & $3d^{10}4s^24p^6$ & d-shell adds 5 Hopf modes \\
Xe & 54 & $4d^{10}5s^25p^6$ & Complete period 5 \\
Rn & 86 & $4f^{14}..6p^6$ & f-shell: 7 highest modes \\
\bottomrule
\end{tabular}
\end{center}

\section{The Periodic Table as Hopf Atlas}

Each period corresponds to completing a Hopf multiplet at level $n$, 
with capacity $2n^2$. Period lengths 2, 8, 8, 18, 18, 32, 32 
follow directly. The periodic table is an atlas of Hopf multiplet 
completions, ordered by $Z$.

\section{Why Chemistry Exists}

Chemistry requires atoms with open Hopf channels: incomplete 
outer-shell multiplets that can accept or donate electrons to 
achieve topological closure. Every chemical bond is a Hopf channel 
negotiation. Covalent bonds share Hopf modes; ionic bonds transfer 
them. The driving force: topological closure.

Standard chemistry says atoms react to achieve full outer shells. 
BCT says atoms react to achieve closed Hopf multiplets. BCT is 
the reason why the first statement is true.

\section{Conclusion}

Atomic structure, the Pauli principle, the octet rule, and noble 
gas inertness all derive from BCT vacuum topology. The periodic 
table is a Hopf fibration atlas. Every noble gas is a topological 
completion. Every chemical bond is a Hopf channel negotiation.

\begin{acknowledgments}
The instinct that noble gases must be topologically complete 
structures --- not merely energetically stable --- was stated 
by Michel Cabri\'e before the calculation was done. 
The instinct was correct. The universe agreed.
\end{acknowledgments}

\begin{thebibliography}{99}
\bibitem{AppJH} M.~R.~Cabri\'e, \textit{Appendix JH: The Hopfion 
Interior and OHC}, Zenodo (2026), doi:10.5281/zenodo.18975018

\bibitem{L1} M.~R.~Cabri\'e, \textit{BCT Letter 1}, Zenodo (2026), 
doi:10.5281/zenodo.18958207

\bibitem{L71} M.~R.~Cabri\'e, \textit{BCT Letter 71: Baryogenesis}, 
Zenodo (2026).

\bibitem{White2026} H.~White, J.~Vera, A.~Sylvester, L.~Dudzinski, 
\textit{Emergent quantization from a dynamic vacuum}, 
Phys.\ Rev.\ Research \textbf{8}, 013264 (2026).
\end{thebibliography}

\end{document}

% (Standalone file: BCT_Letter78_AtomicStructure.tex)

The OHC Bessel resonance modes generate electron shells with
capacity $2n^2$ as a derived theorem. Shell filling follows
the Janet (left-step) sequence. The Bohr radius is derived
to sub-0.1\%. This letter establishes the BCT foundation
for atomic structure theory.

\textit{See standalone file: BCT\_Letter78\_AtomicStructure.tex/.pdf}

% ── LETTER 79 ────────────────────────────────────────────────
\newpage
\section*{Letter 79: Morphic Resonance and the OHC Condensate}
\addcontentsline{toc}{section}{Letter 79: Morphic Resonance}
% Written in response to enquiry from Prof.\ Rupert Sheldrake.

The OHC condensate provides a physical substrate for
morphic-resonance-type phenomena: two memory channels
($S_\mathrm{Hopf} = 5329$ local; $v_\mathrm{ph} = 6.78c$ global).

\textit{See standalone file: BCT\_Letter79\_MorphicResonance.tex/.pdf}

% ── LETTER 80 ────────────────────────────────────────────────
\newpage
\section*{Letter 80: Molecular Geometry from Hopf Topology}
\addcontentsline{toc}{section}{Letter 80: Molecular Geometry}

The universal BCT Bond Angle Theorem:
\begin{equation}
  \cos\theta = -\frac{1 - n_\mathrm{lp}/8}{3}
\end{equation}
yields H$_2$O at $104.478^\circ$ (error $-0.0001\%$) and
CH$_4$ at $109.471^\circ$ (exact). The algebraic identity
$r_\mathrm{tet}(1+r_\mathrm{tet}) = 1/8$ is proven.

\textit{See standalone file: BCT\_Letter80\_MolecularGeometry.tex/.pdf}

% ── LETTER 81 ────────────────────────────────────────────────
\newpage
\section*{Letter 81: The Two Fundamental Theorems of BCT Chemistry}
\addcontentsline{toc}{section}{Letter 81: Two Fundamental Theorems}

\textbf{Theorem 1 (Pauli Exclusion):} $|H| \leq n$ per OHC mode,
derived from $H \in \mathbb{Z}$ and Josephson--Bessel capacity.
\textbf{Theorem 2 (Noble Gas Inertness):} closed Hopf multiplets
have no vacant acceptor modes; bonding is topologically impossible.

\textit{See standalone file: BCT\_Letter81\_TwoTheorems.tex/.pdf}

% ── LETTER 82 ────────────────────────────────────────────────
\newpage
\section*{Letter 82: Carbon: The Topological Basis of Life}
\addcontentsline{toc}{section}{Letter 82: Carbon}

$S_H(\mathrm{C}) = 1.00$ is unique in the $n=2$ shell.
The Life Axis: H $\to$ C $\to$ Si $\to$ Ge $\to$ Sn $\to$
Pb $\to$ Fl, all with $S_H = 1$. Carbon's capacity for
self-bonding chains is a topological necessity, not an
energetic accident.

\textit{See standalone file: BCT\_Letter82\_Carbon.tex/.pdf}

% ── LETTER 83 ────────────────────────────────────────────────
\newpage
\section*{Letter 83: Transition Metal Anomalies as BCT Hopf Symmetry Points}
\addcontentsline{toc}{section}{Letter 83: Transition Metal Anomalies}

Anomalous configurations of Cr, Mo (d half-fill attractor,
$S_H(d) = 0.50$), Cu, Ag, Au (d-closure attractor, $d^{10}$),
Eu, Am (f half-fill), and Yb, No (f-closure) are all derived
consequences of the Hopf symmetry attractor theorem.
Gold's 521\,nm absorption (yellow colour) is derived.
Mercury's room-temperature liquidity follows from double
Hopf closure.

\textit{See standalone file: BCT\_Letter83\_TransitionMetals.tex/.pdf}

% ── LETTER 84 ────────────────────────────────────────────────
\newpage
\section*{Letter 84: The BCT Periodic Table}
\addcontentsline{toc}{section}{Letter 84: BCT Periodic Table}

A complete reclassification of all 118 elements by Hopf
multiplet topology into 11 classes. Two novel layouts:
the Janet left-step arrangement and the circular
\textit{In the Round} arrangement. Three topological axes
are identified and proven: the Life Axis ($S_H = 1$,
6\,o'clock), the Noble Metal Axis (d-closure, 7\,o'clock),
and the d/f Half-Fill Axis (Hopf attractor, 4\,o'clock).

\textit{See standalone file: BCT\_Letter84\_PeriodicTable.tex/.pdf}

% ── LETTER 85 ────────────────────────────────────────────────
\newpage
\section*{Letter 85: Electronegativity from Void Geometry}
\addcontentsline{toc}{section}{Letter 85: Electronegativity}

BCT screening constants from void geometry:
$\sigma_\mathrm{inner} = r_\mathrm{tet}/r_\mathrm{oct} = 0.5426$,
$\sigma_\mathrm{same} = r_\mathrm{tet}/(1+r_\mathrm{tet}) = 0.1010$.
Period 2 elements reproduced to sub-3\%; O at $+0.2\%$,
P at $+0.1\%$. The Polar Bond Angle Theorem corrects NH$_3$
from $-0.78\%$ to $-0.027\%$.

\textit{See standalone file: BCT\_Letter85\_Electronegativity.tex/.pdf}

% ════════════════════════════════════════════════════════════
% PART II: APPENDICES
% ════════════════════════════════════════════════════════════

\newpage
\part*{Part II: Appendices KA1--KC2}

% ── APPENDIX KA1 ─────────────────────────────────────────────
\newpage
\section*{Appendix KA1: OHC Mode Structure and Shell Capacity}
\addcontentsline{toc}{section}{Appendix KA1: OHC Mode Structure}
\textit{See standalone file: BCT\_Appendix\_KA1.tex/.pdf}

% ── APPENDIX KA2 ─────────────────────────────────────────────
\newpage
\section*{Appendix KA2: OHC Phase Memory and the Repetition Theorem}
\addcontentsline{toc}{section}{Appendix KA2: Phase Memory}
\textit{See standalone file: BCT\_Appendix\_KA2.tex/.pdf}

% ── APPENDIX KA3 ─────────────────────────────────────────────
\newpage
\section*{Appendix KA3: The $r_\mathrm{tet}$ Identity and Bond Angle Series}
\addcontentsline{toc}{section}{Appendix KA3: r-tet Identity}
\textit{See standalone file: BCT\_Appendix\_KA3.tex/.pdf}

% ── APPENDIX KA4 ─────────────────────────────────────────────
\newpage
\section*{Appendix KA4: Formal $|H| \leq n$ Proof}
\addcontentsline{toc}{section}{Appendix KA4: Pauli Proof}
\textit{See standalone file: BCT\_Appendix\_KA4.tex/.pdf}

% ── APPENDIX KB1 ─────────────────────────────────────────────
\newpage
\section*{Appendix KB1: Carbon $S_H$ Index and Bond Compression}
\addcontentsline{toc}{section}{Appendix KB1: Carbon S-H Index}
\textit{See standalone file: BCT\_Appendix\_KB1.tex/.pdf}

% ── APPENDIX KB2 ─────────────────────────────────────────────
\newpage
\section*{Appendix KB2: d-Shell Hopf Symmetry}
\addcontentsline{toc}{section}{Appendix KB2: d-Shell Hopf Symmetry}
\textit{See standalone file: BCT\_Appendix\_KB2.tex/.pdf}

% ── APPENDIX KB3 ─────────────────────────────────────────────
\newpage
\section*{Appendix KB3: f-Shell Hopf Symmetry}
\addcontentsline{toc}{section}{Appendix KB3: f-Shell Hopf Symmetry}
\textit{See standalone file: BCT\_Appendix\_KB3.tex/.pdf}

% ── APPENDIX KC1 ─────────────────────────────────────────────
\newpage
\section*{Appendix KC1: BCT Periodic Table --- Full Classification}
\addcontentsline{toc}{section}{Appendix KC1: Periodic Table Classification}
\textit{See standalone file: BCT\_Appendix\_KC1.tex/.pdf}

% ── APPENDIX KC2 ─────────────────────────────────────────────
\newpage
\section*{Appendix KC2: The Three Topological Axes}
\addcontentsline{toc}{section}{Appendix KC2: Three Topological Axes}
\textit{See standalone file: BCT\_Appendix\_KC2.tex/.pdf}

% ── CITATION BLOCK ───────────────────────────────────────────
\newpage
\section*{How to Cite This Volume}

To cite the complete Chemistry Series volume:
\begin{quote}
Cabri\'{e}, M.R. (2026). \textit{BCT Chemistry Series,
Volume K: Letters 78--85 and Appendices KA1--KC2}.
Zenodo. \texttt{doi:10.5281/zenodo.XXXXXXX}
\end{quote}

To cite individual letters, use the standalone PDFs/DOIs.
Individual \texttt{.tex} files for all letters and appendices
are provided alongside this volume for Overleaf compilation.

\section*{Acknowledgements}

CSSC (Claude The Super Super Computer, named by Jeff) assisted
with document preparation, numerical verification, and
programme management throughout the BCT Chemistry Series.
Nyssa and Tegan provided supervisory oversight.

\end{document}
